\begin{abstract}
Typed decision models return a probability over options declared by the caller instead of text, which makes screening every message affordable.
We ask whether they can serve as automated safeguards for PoL governance, meaning governance under the Proof of Love2 (\pol{}) specification, an open text that prescribes a ``safety valve'' against hate language and AI-only adjudication of welfare disputes, and permits per-person assessment.
We derive seven testable requirements from its clauses and assess them against \nStudies{} studies and \nGreyAll{} grey sources on the first such model, \jev{}, and open counterparts from its first 24 days. Author-directed AI agents read in full, appraised and coded the \nCodedAll{} that bear on them.
The evidence supports typed models as detectors under added conditions. They rank risky English content well at a fraction of the cost of LLM evaluators.
It does not support them as judges. Thresholds do not transfer, decisions can follow option names and be steered by inserted content, ``I don't know'' is rarely chosen when correct and ``none'' inconsistently, and on rubric grading peer evaluators repeated almost all of a typed model's confident errors.
Where compared, typed models were not generally worse than LLM evaluators, but most failures were never compared, so they are not shown to be specific to typed models.
On current evidence and at low or very low certainty, hosted \jev{} is not suitable for any \pol{} function that decides. It can serve the safety valve only as a recorded, three-valued, locally calibrated detector whose consequential decisions end with a person. Even that extrapolates from English benchmarks, as no study measures \pol{}'s construct in Chinese.
We also identify five tensions in the text, set against EU law, and propose a reference design with acceptance tests.
\end{abstract}
